\section{What the measurements predict}
\label{sec:roadmap}

Each planned feature can now be given a measured expectation before it is
built (Table~\ref{tab:roadmap}), so each one becomes a test of the model as
well as an improvement.

\begin{table}[ht]
\centering\small
\begin{tabular}{p{0.24\linewidth}p{0.68\linewidth}}
\toprule
next step & what the measurements predict \\
\midrule
restarts & less sensitivity to variable order and signs (the twofold shift
           should shrink toward MiniSat's 1.1); not necessarily fewer
           conflicts on pigeonhole \\
clause minimization & learnt clauses about 19\% shorter and 13--27\% fewer
           conflicts (Section~\ref{sec:minisat}) \\
blocking literals & saves at most the cost of skips: 8--15\% of cycles on
           pigeonhole, 4--19\% on random 3-SAT, up to 87\% on fullN; more once
           data leaves L2 \\
cheaper search & the largest target, 62--82\% of cycles on pigeonhole and
           random 3-SAT; resuming the scan where the last one
           stopped~\cite{gent2013watched} would shorten it \\
VSIDS heap & no measurable effect at 250 variables; matters at $10^5$ \\
arena storage & removes one pointer indirection per visit, the cost measured
           in Section~\ref{sec:memory} \\
binary DRAT & smaller proofs, less of the
           \numProofOverheadMin--\numProofOverheadMax\% logging overhead \\
\bottomrule
\end{tabular}
\caption{Measured expectations for the planned features.}
\label{tab:roadmap}
\end{table}
